\section{Determination of the process-zone parameters}
\label{sec:derivation}

\subsection{Mellin transform and functional equation}
\label{subsec:mellin}

To construct the local solution, we introduce the dimensionless radial
coordinate
\begin{equation*}
  \rho=\frac{r}{d_i}
\end{equation*}
and apply the Mellin integral transform to the boundary-value problem
\eqref{eq:interface-boundary-conditions}--\eqref{eq:matched-far-field}.
Let $p$ be the complex Mellin variable. The transforms of the correction normal stress outside the zone and the
normal displacement gradient within the zone are defined, respectively, by
\begin{equation*}
\begin{aligned}
  \Phi_i^+(p)
    &=\int_1^\infty
      \widehat{\sigma}_i(\rho d_i)\rho^p\,\dd\rho,\\
  \Phi_i^-(p)
    &=\frac{E_i}{2(1-\nu_i^2)}
      \int_0^1
      \left.
      \frac{\partial u_\theta}{\partial r}
      \right|_{\substack{r=\rho d_i\\\theta=\beta_i}}
      \rho^p\,\dd\rho.
\end{aligned}
\end{equation*}
Because $\widehat{\sigma}_i(r)=o(r^{-1})$ as $r\to\infty$, the integral
for $\Phi_i^+$ defines an analytic function for
$\operatorname{Re}p<0$. Near the corner, the displacement gradient has
order $r^{\lambda_i}$, where $\lambda_i$ is defined below, so
$\Phi_i^-$ is analytic for
$\operatorname{Re}p>-1-\lambda_i$. The superscripts on $\Phi_i^+$ and
$\Phi_i^-$ identify functions assigned to the left and right half-planes,
respectively; they do not indicate the sign of a physical quantity.

On the active segment $0<\rho<1$, the correction traction is
$\sigma_i-CQ_i d_i^\lambda\rho^\lambda$, and hence
\begin{equation*}
  \int_0^1
  \left(\sigma_i-CQ_i d_i^\lambda\rho^\lambda\right)
  \rho^p\,\dd\rho
  =
  \frac{\sigma_i}{p+1}
  -\frac{CQ_i d_i^\lambda}{p+\lambda+1}.
\end{equation*}
Eliminating the transforms of the interior field yields the functional
equation
\begin{equation}
\label{eq:wiener-hopf-functional}
  \Phi_i^+(p)
  +\frac{\sigma_i}{p+1}
  -\frac{CQ_i d_i^\lambda}{p+\lambda+1}
  =-\tan(\pi p)G_i(p)\Phi_i^-(p),
\end{equation}
which is initially valid in the common convergence strip
\begin{equation*}
  \max\{-1-\lambda,-1-\lambda_i\}
  <\operatorname{Re}p<0.
\end{equation*}
Here, $\operatorname{Re}p$ denotes the real part of $p$. In this strip,
both sides represent the same analytic relation. The functional equation
then provides the analytic (or, in the presence of the explicit rational
terms, meromorphic) continuation of the required combinations to the
factorization contour $\operatorname{Re}p=0$. No convergence of the
integral defining $\Phi_i^+$ on the imaginary axis is asserted. The
kernel $G_i(p)$ is
\begin{equation*}
  G_i(p)=\frac{D_i(p)\cos(\pi p)}{D(p)\sin(\pi p)}.
\end{equation*}
Here, $D(p)$ is the characteristic determinant for the corner with an
interfacial shear crack, given in Appendix~\ref{app:asymptotics}.
For plane strain, define the dimensionless elastic ratio $e$ and the
material parameters $\kappa_i$ by
\begin{equation*}
  e=\frac{1+\nu_2}{1+\nu_1}\frac{E_1}{E_2},
  \qquad
  \kappa_i=3-4\nu_i.
\end{equation*}
The characteristic determinants $D_1(p)$ and $D_2(p)$ for configurations
with a process zone in materials~1 and~2, respectively, are
\begin{equation*}
\begin{aligned}
  D_1(p)={}&2e(1+\kappa_2)
    [\cos(2p\alpha)-\cos(2\alpha)]
    \{\sin^2[p(\pi-\alpha)]-p^2\sin^2\alpha\}\\
  &+(1+\kappa_1)
    [\sin(2p\alpha)+p\sin(2\alpha)]
    [\sin(2p(\pi-\alpha))-p\sin(2\alpha)],\\[2pt]
  D_2(p)={}&2(1+\kappa_1)
    [\cos(2p(\pi-\alpha))-\cos(2\alpha)]
    [\sin^2(p\alpha)-p^2\sin^2\alpha]\\
  &+e(1+\kappa_2)
    [\sin(2p(\pi-\alpha))-p\sin(2\alpha)]
    [\sin(2p\alpha)+p\sin(2\alpha)].
\end{aligned}
\end{equation*}
In particular, the corner singularity exponent $\lambda_i$ in the
presence of a zone in material~$i$ is determined by
\begin{equation*}
  D_i(-1-\lambda_i)=0,
\end{equation*}
where $\lambda_i$ is the smallest real root in $(-1,0)$ continuously
connected to the corresponding physical branch. As before,
$\alpha=\pi/2$ is understood only as a degenerate limit.

\subsection{Wiener--Hopf factorization}
\label{subsec:wiener-hopf}

We solve Eq.~\eqref{eq:wiener-hopf-functional} by the Wiener--Hopf
method~\cite{noble1958wienerhopf}, following the approach used for a related
boundary-value problem~\cite{KaminskyDudykPolishchuk2024}.
The factors $G_i^\pm(p)$ are defined by the Cauchy integral
\begin{equation*}
  \exp\left\{
    \frac{1}{2\pi\ii}
    \int_{-\ii\infty}^{\ii\infty}
      \frac{\log G_i(z)}{z-p}\,\dd z
  \right\}
  =
  \begin{cases}
    G_i^+(p), & \operatorname{Re}p<0,\\
    G_i^-(p), & \operatorname{Re}p>0,
  \end{cases}
\end{equation*}
where $z$ is the complex integration variable and $\ii^2=-1$.
The contour runs upward along the imaginary axis. We also introduce
the auxiliary gamma-function factors
\begin{equation*}
  K^\pm(p)=
  \frac{\Gamma(1\mp p)}{\Gamma(\tfrac12\mp p)}.
\end{equation*}
where $\Gamma$ is the Euler gamma function.
With the normalization $G_i^\pm(p)\to1$ at infinity, the solution of
the functional equation is
\begin{equation}
\label{eq:wiener-hopf-solution}
\begin{aligned}
  \Phi_i^+(p)={}&
  -\frac{pG_i^+(p)}{K^+(p)}
  \Bigg\{
    \frac{\sigma_i}{p+1}
    \left[
      \frac{K^+(p)}{pG_i^+(p)}
      +\frac{K^+(-1)}{G_i^+(-1)}
    \right]\\
  &\hspace{37mm}
    -\frac{CQ_i d_i^\lambda}{p+\lambda+1}
    \left[
      \frac{K^+(p)}{pG_i^+(p)}
      +\frac{K^+(-\lambda-1)}
      {(\lambda+1)G_i^+(-\lambda-1)}
    \right]
  \Bigg\},
  \qquad \operatorname{Re}p<0,\\[3pt]
  \Phi_i^-(p)={}&
  K^-(p)G_i^-(p)
  \Bigg\{
    \frac{\sigma_iK^+(-1)}{(p+1)G_i^+(-1)}
    -\frac{CQ_i d_i^\lambda K^+(-\lambda-1)}
    {(\lambda+1)(p+\lambda+1)G_i^+(-\lambda-1)}
  \Bigg\},
  \qquad \operatorname{Re}p>0.
\end{aligned}
\end{equation}
For the kernels used below, $G_i(\ii t)$ is positive and even in the
real variable $t$, and tends to unity as $|t|\to\infty$.
Thus, for real $p<0$, the plus factor can be evaluated numerically
using the real integral
\begin{equation*}
  \log G_i^+(p)
  =-\frac{p}{\pi}\int_0^\infty
    \frac{\log G_i(\ii t)}{t^2+p^2}\,\dd t,
  \qquad p<0.
\end{equation*}
In the numerical calculations, we check that the kernel is real and
positive on the factorization contour and approaches unity at its
far end.

\subsection{Process-zone length}
\label{subsec:zone-length}

As $p\to+\infty$, the second relation in
Eq.~\eqref{eq:wiener-hopf-solution} gives
\begin{equation*}
  \Phi_i^-(p)\sim\frac{1}{\sqrt{p}}
  \left[
    \frac{\sigma_iK^+(-1)}{G_i^+(-1)}
    -\frac{CQ_i d_i^\lambda K^+(-\lambda-1)}
    {(\lambda+1)G_i^+(-\lambda-1)}
  \right].
\end{equation*}
Applying an Abelian-type theorem~\cite{noble1958wienerhopf} to the
tip asymptotics~\eqref{eq:zone-tip-asymptotics} also gives
\begin{equation*}
  \Phi_i^-(p)\sim-\frac{k_i}{\sqrt{2pd_i}}.
\end{equation*}
Comparing these expressions yields the local stress intensity factor
\begin{equation*}
  k_i=-\sqrt{2d_i}
  \left[
    \frac{\sigma_iK^+(-1)}{G_i^+(-1)}
    -\frac{CQ_i d_i^\lambda K^+(-\lambda-1)}
    {(\lambda+1)G_i^+(-\lambda-1)}
  \right].
\end{equation*}
Imposing $k_i=0$, which removes the inverse-square-root singular term
from the assumed tip asymptotics, gives the zone length:
\begin{equation}
\label{eq:process-zone-length}
  d_i=
  \left[
    \frac{CQ_iK^+(-\lambda-1)G_i^+(-1)}
    {\sigma_i(1+\lambda)K^+(-1)G_i^+(-\lambda-1)}
  \right]^{-1/\lambda}.
\end{equation}
Since $-1<\lambda<0$, a real, positive $d_i$ requires $CQ_i>0$,
consistent with the physical admissibility condition in
Section~\ref{subsec:boundary-conditions}. With geometry and material
properties fixed, Eq.~\eqref{eq:process-zone-length} also gives the
power law $d_i\propto C^{-1/\lambda}$ on the admissible loading branch.

Alternatively, with the geometry, interfacial crack length, elastic
properties, and applied loading fixed, the same equation gives the
dependence on the normal cohesive traction:
\begin{equation*}
  d_i\propto \sigma_i^{1/\lambda}
  =\sigma_i^{-1/|\lambda|}.
\end{equation*}
Because $-1<\lambda<0$, decreasing $\sigma_i$ increases the
process-zone length.

Using the relation
\begin{equation*}
  Q_i=g_1q_i l^{\lambda_0-\lambda},
\end{equation*}
where $g_1$ is the intact-corner angular coefficient and $q_i$ is a
dimensionless amplitude coefficient, both given in
Appendix~\ref{app:asymptotics}, we obtain from
Eq.~\eqref{eq:process-zone-length}
\begin{equation}
\label{eq:normalized-process-zone-length}
  \frac{d_i}{l}=
  \left[
    \frac{Cl^{\lambda_0}}{\sigma_i}
    \frac{g_1q_iK^+(-1-\lambda)G_i^+(-1)}
    {(1+\lambda)K^+(-1)G_i^+(-1-\lambda)}
  \right]^{-1/\lambda}.
\end{equation}
This relation also provides a check on the small-scale assumption
$d_i/l\ll1$ used in the local formulation.

\subsection{Opening and energy parameter}
\label{subsec:zone-parameters}

Define the separation of the zone faces at the corner by
$\delta_i=\jump{u_\theta(0,\beta_i)}$.
The definition of $\Phi_i^-(p)$ gives
\begin{equation*}
  \delta_i
  =-2
    \left.
      \left(\frac{\partial u_\theta}{\partial r}\right)^*
    \right|_{p=0,\,\theta=\beta_i}
  =-\frac{4(1-\nu_i^2)d_i}{E_i}\Phi_i^-(0).
\end{equation*}
Here, the asterisk denotes the dimensional transform moment:
for $f(r)=\partial u_\theta(r,\beta_i)/\partial r$,
$f^*(p)=\int_0^{d_i}f(r)(r/d_i)^p\,\dd r$.
The factor of two accounts for the symmetric displacement of the two
zone faces. Substitution of Eq.~\eqref{eq:wiener-hopf-solution},
together with Eq.~\eqref{eq:process-zone-length}, gives
\begin{equation}
\label{eq:process-zone-opening}
  \delta_i=
  -\frac{4(1-\nu_i^2)\sigma_i\lambda K^+(-1)}
  {E_i\sqrt{\pi G_i(0)}(1+\lambda)G_i^+(-1)}d_i.
\end{equation}
For fixed material properties and geometry, the coefficient multiplying
$d_i$ in Eq.~\eqref{eq:process-zone-opening} is constant. The opening
therefore scales with the zone length as the applied loading varies:
$\delta_i\propto d_i$. If the applied loading is fixed and $\sigma_i$
is varied, the preceding power law instead gives
\begin{equation*}
  \delta_i\propto \sigma_i d_i
  \propto \sigma_i^{1+1/\lambda}
  =\sigma_i^{1-1/|\lambda|}.
\end{equation*}
The exponent $1-1/|\lambda|$ is negative. A decrease in the resistance
to separation therefore increases both the zone length and its opening.

Within the adopted model, the work $W_i$ done against the resistance
to separation, per unit out-of-plane thickness, is defined
by~\cite{Cherepanov1986}
\begin{equation*}
  W_i=\sigma_i\int_0^{d_i}
      \jump{u_\theta(r,\beta_i)}\,\dd r
  =-\frac{4(1-\nu_i^2)\sigma_i}{E_i}d_i^2\Phi_i^-(1).
\end{equation*}
Using Eq.~\eqref{eq:wiener-hopf-solution}, we obtain
\begin{equation*}
  W_i=-\frac{8\lambda(1-\nu_i^2)\sigma_i^2}
  {\pi(\lambda+2)[G_i^+(-1)]^2E_i}d_i^2.
\end{equation*}
Differentiating with respect to the zone length along the family of
equilibrium states, with the geometry and material properties fixed,
defines the zone-related energy parameter $J_i$:
\begin{equation}
\label{eq:process-zone-energy-release-rate}
  J_i=\frac{\dd W_i}{\dd d_i}
  =-\frac{16\lambda(1-\nu_i^2)\sigma_i^2}
  {\pi(\lambda+2)[G_i^+(-1)]^2E_i}d_i.
\end{equation}
For the numerical analysis, introduce the dimensionless loading,
opening, and energy parameters $\sigmap$, $\delta_i'$, and $J_i'$,
respectively:
\begin{equation*}
  \sigmap=\frac{Cl^{\lambda_0}}{\sigma_i},
  \qquad
  \delta_i'=\frac{E_i\delta_i}
  {4(1-\nu_i^2)\sigma_i l},
  \qquad
  J_i'=\frac{E_iJ_i}
  {4(1-\nu_i^2)\sigma_i^2l}.
\end{equation*}
The primes here denote normalization, not differentiation.
Equations~\eqref{eq:normalized-process-zone-length},
\eqref{eq:process-zone-opening}, and
\eqref{eq:process-zone-energy-release-rate} then take the form
\begin{equation*}
\begin{aligned}
  \frac{d_i}{l}
    &=\left[
      \sigmap\,
      \frac{g_1q_iK^+(-1-\lambda)G_i^+(-1)}
      {(1+\lambda)K^+(-1)G_i^+(-1-\lambda)}
    \right]^{-1/\lambda},\\[2pt]
  \delta_i'
    &=-\frac{\lambda K^+(-1)}
      {\sqrt{\pi G_i(0)}(1+\lambda)G_i^+(-1)}
      \frac{d_i}{l},\\[2pt]
  J_i'
    &=-\frac{4\lambda}
      {\pi(\lambda+2)[G_i^+(-1)]^2}
      \frac{d_i}{l}.
\end{aligned}
\end{equation*}
All three relations use the corner exponent $\lambda$ for the
interfacial shear crack, i.e., the root of $D(-1-\lambda)=0$.
The exponents $\lambda_1$ and $\lambda_2$, which describe the corner
singularity in the presence of a zone, do not replace $\lambda$
in the formulas for the zone parameters.

Equating Eq.~\eqref{eq:process-zone-energy-release-rate} to an
independently specified critical value $J_{i\mathrm{c}}$ for material~$i$
provides a condition for the applied loading at which the process zone
develops into a Mode~I crack branching from the original interfacial
crack. This criterion concerns local branching initiation within the
adopted model.
